\documentclass[11pt]{article}
\usepackage[margin=1in]{geometry}
\usepackage{amsmath,amssymb,amsthm}
\usepackage{hyperref}
\usepackage{enumitem}

\newtheorem{theorem}{Theorem}
\newtheorem{lemma}[theorem]{Lemma}
\newtheorem{corollary}[theorem]{Corollary}
\theoremstyle{definition}
\newtheorem{definition}[theorem]{Definition}
\theoremstyle{remark}
\newtheorem{remark}[theorem]{Remark}

\usepackage{xurl}
\let\oldus\_
\renewcommand{\_}{\oldus\allowbreak}

\newcommand{\Q}{\mathbb{Q}}
\newcommand{\Qb}{\overline{\mathbb{Q}}}
\newcommand{\GQ}{G_{\mathbb{Q}}}

\title{A disproof of Conjecture 00000002416\\[2pt]
\large No open subgroup of $\mathrm{Gal}(\overline{\mathbb{Q}}/\mathbb{Q})$ is finitely generated}
\date{}

\begin{document}
\sloppy
\maketitle

\section{The conjecture}

Conjecture 00000002416 reads (English text; the Chinese text says the same):
\begin{quote}
Conjecture: The spectrum of the minimal number of generators of open subgroups of $G_\Q$ is $\{2,3\}$;
the density of $d=3$ subgroups is $0$ (an explicit density -- almost all subgroups are 2-generated);
the Shafarevich spectrum is complete.
\end{quote}
The statement is a conjunction of three claims separated by semicolons. We refute the first one, which
the second one also presupposes (it speaks of the ``$d=3$ subgroups''). Hence the conjecture is false.

\section{Reading and definitions}

\begin{itemize}[leftmargin=*]
\item $\GQ=\mathrm{Gal}(\Qb/\Q)$ is the absolute Galois group of $\Q$, with its Krull topology. In Lean
  this is Mathlib's \texttt{Field.absoluteGaloisGroup Q} (with Lean's $\Q$ written as Q here), the group of $\Q$-algebra automorphisms of
  \texttt{AlgebraicClosure Q} with the Krull topology.
\item A subset $S$ of a topological group $H$ \emph{topologically generates} $H$ if the subgroup
  $\langle S\rangle$ generated by $S$ is dense in $H$ (Lean: \texttt{TopGenerates}).
\item The \emph{minimal number of generators} of $H$ is
  $d(H)=\inf\{|S| : S\subseteq H \text{ finite}, S \text{ topologically generates } H\}\in\mathbb{N}\cup\{\infty\}$,
  with $d(H)=\infty$ (Lean: $\top$ in \texttt{ENat}) when no finite set topologically generates $H$
  (Lean: \texttt{minGen}).
\item An open subgroup $U\le\GQ$ carries the subspace topology. The \emph{spectrum} is
  $\{d(U): U\le \GQ \text{ an open subgroup}\}\subseteq \mathbb{N}\cup\{\infty\}$ (Lean: \texttt{genSpectrum}).
  It contains $d(\GQ)$, since $\GQ$ is an open subgroup of itself.
\item We also treat abstract generation: $d_{\mathrm{abs}}(H)$ is the least size of a finite $S$ with
  $\langle S\rangle=H$ (Lean: \texttt{minGenAbs}, spectrum \texttt{genSpectrumAbs}).
\end{itemize}

\section{The result}

\begin{theorem}\label{thm:main}
No open subgroup of $\GQ$ is topologically finitely generated. Hence $d(U)=\infty$ for every open subgroup $U$,
the spectrum is $\{\infty\}$, and in particular it is not $\{2,3\}$; neither $2$ nor $3$ belongs to it. The same
holds for abstract generation.
\end{theorem}

\subsection*{Step 1: quadratic characters}
For a prime $p$ fix $\sqrt p\in\Qb$ with $(\sqrt p)^2=p$ (any choice). For $\sigma\in\GQ$ we have
$\sigma(\sqrt p)^2=\sigma(p)=p$, so $\sigma(\sqrt p)=\pm\sqrt p$. Put $\chi_p(\sigma)=1$ if
$\sigma(\sqrt p)=\sqrt p$ and $\chi_p(\sigma)=-1$ otherwise. Because $\sqrt p\neq -\sqrt p$ (characteristic $0$,
$p\neq0$), a case check shows that $\chi_p:\GQ\to\{\pm1\}=\mathbb{Z}^\times$ is a group homomorphism.

Its kernel contains $\mathrm{Gal}(\Qb/\Q(\sqrt p))$, the fixing subgroup of the finite extension
$\Q(\sqrt p)$, which is open in the Krull topology. A subgroup containing an open subgroup is open, so
$\ker\chi_p$ is open.

\subsection*{Step 2: distinct primes give distinct characters}
Let $p\neq q$ be primes with $\chi_p=\chi_q$. For every $\sigma$,
$\sigma(\sqrt p\sqrt q)=\chi_p(\sigma)\chi_q(\sigma)\sqrt p\sqrt q=\sqrt p\sqrt q$. The extension $\Qb/\Q$ is
Galois, so by the Galois correspondence for (possibly infinite) Galois extensions the elements fixed by all
of $\GQ$ form $\Q$. Hence $\sqrt p\sqrt q=r\in\Q$ and $r^2=pq$. A natural number that is a square in $\Q$ is a
square in $\mathbb{N}$, so $pq=m^2$. Then $p\mid m$, say $m=pk$, so $q=pk^2$ and $p\mid q$, forcing $p=q$, a
contradiction.

\subsection*{Step 3: counting}
Let $f,g:H\to M$ be homomorphisms of groups, $H$ a topological group, with open kernels, and suppose $f=g$
on a set $S$ that topologically generates $H$. The equalizer $E=\{x: f(x)=g(x)\}$ is a subgroup containing the
open subgroup $\ker f\cap\ker g$, so $E$ is open, hence closed. It contains $\langle S\rangle$, so it contains
its closure, which is $H$. Thus $f=g$.

Suppose a finite set $T$ topologically generates $\GQ$. The map sending a prime $p$ to the restriction
$\chi_p|_T\in\{\pm1\}^T$ is injective by Step 3 and Step 2. The primes are infinite and $\{\pm1\}^T$ is finite,
a contradiction. So $\GQ$ is not topologically finitely generated.

\subsection*{Step 4: open subgroups}
$\GQ$ is compact. Let $U$ be an open subgroup and suppose a finite $S\subseteq U$ topologically generates $U$ in
the subspace topology. The cosets $gU$ form an open cover of $\GQ$, so finitely many, $r_1U,\dots,r_kU$, cover it.
Let $H=\langle S\cup\{r_1,\dots,r_k\}\rangle\le\GQ$. Since the inclusion $U\to\GQ$ is continuous, $U$ lies in the
closure $\overline H$; the $r_i$ lie in $H$; and $\overline H$ is a subgroup, so it contains every $r_iu$ and
hence equals $\GQ$. So the finite set $S\cup\{r_i\}$ topologically generates $\GQ$, contradicting Step 3.

Finally, $\langle S\rangle=U$ implies that $\langle S\rangle$ is dense in $U$, so the abstract version follows.
Since $\GQ$ itself is open, the spectrum is exactly $\{\infty\}$. \qed

\section{Lean formalization}

The Lean file \texttt{Conjecture2416/Basic.lean} (296 lines, only import \texttt{Mathlib}, namespace
\texttt{C2416}) contains:
\begin{itemize}[leftmargin=*]
\item \texttt{topFG\_of\_open}: Step 4 for any compact topological group.
\item \texttt{eq\_of\_eqOn\_gen}: Step 3, the equalizer argument.
\item \texttt{sqrtQ}, \texttt{chi}, \texttt{chiHom}, \texttt{chiHom\_ker\_isOpen}: Step 1.
\item \texttt{mem\_range\_of\_fixed} (from Mathlib's \texttt{InfiniteGalois.fixedField\_fixingSubgroup}),
  \texttt{not\_isSquare\_mul}, \texttt{chiHom\_injective}: Step 2.
\item \texttt{gal\_not\_topFG}, \texttt{absGal\_not\_topFG}, \texttt{open\_subgroup\_not\_topFG},
  \texttt{open\_subgroup\_not\_fg}, \texttt{minGen\_open\_subgroup}, \texttt{genSpectrum\_absGal},
  \texttt{genSpectrumAbs\_absGal}.
\end{itemize}
The main theorem is
\begin{verbatim}
theorem conjecture_2416_false :
    genSpectrum (Field.absoluteGaloisGroup Q) != {2, 3} /\
    2 not-in genSpectrum (Field.absoluteGaloisGroup Q) /\
    3 not-in genSpectrum (Field.absoluteGaloisGroup Q)
\end{verbatim}
(with the Lean symbols for not-equal, and, not-in written in ASCII), together with \texttt{conjecture\_2416\_false\_abstract}, the same
statement for \texttt{genSpectrumAbs}, and \texttt{genSpectrum\_absGal : genSpectrum (Field.absoluteGaloisGroup Q) = \{top\}}.

One technical point: \texttt{AlgebraicClosure Q} carries two $\Q$-algebra structures that are equal but not
reducibly so (\texttt{AlgebraicClosure.instAlgebra}, used in the definition of
\texttt{Field.absoluteGaloisGroup}, and \texttt{DivisionRing.toRatAlgebra}). The file gives the former local
priority, so all statements refer to the same automorphism group as \texttt{Field.absoluteGaloisGroup Q}.

\paragraph{Axiom audit.} \texttt{lake build} succeeds. \texttt{\#print axioms} for
\texttt{conjecture\_2416\_false}, \texttt{conjecture\_2416\_false\_abstract}, \texttt{genSpectrum\_absGal} and
\texttt{open\_subgroup\_not\_topFG} reports only \texttt{propext}, \texttt{Classical.choice}, \texttt{Quot.sound}.
There is no \texttt{sorry} and no \texttt{native\_decide}.

\section{Scope}
We refute the first clause for $\GQ=\mathrm{Gal}(\Qb/\Q)$, with topological generation (the standard notion for
profinite groups) and with abstract generation. We do not address the other two clauses, which are not needed:
the conjecture is their conjunction with the first one.

\end{document}
